% Joint follow-up to the IWAI 2026 paper (Hallgren, Hyland, Albarracin).
% Scope agreed on the 1 Oct 2026 call: pluralist agents as the base; the question is
% whether pluralist communities resist two-channel precision collapse better than
% monist ones; belief utility and the motivated cost of mind-change come back in.
% Trust split / selective sharing and motivated ignorance are follow-up papers.
% Target: a Royal Society journal (rolling). Swap in the journal class once chosen.
% Results sections are placeholders until P1/P2 are reviewed.

\documentclass[11pt]{article}
\usepackage[margin=1in]{geometry}
\usepackage{amsmath,amssymb}
\usepackage{graphicx}
\usepackage{booktabs}
\usepackage[numbers]{natbib}
\usepackage{hyperref}

\newcommand{\todo}[1]{\textcolor{red}{[TODO: #1]}}
\usepackage{xcolor}

\title{When Do Pluralist Communities Lock In?\\
\large Precision collapse on the world and social channels}
% Author order to settle with Jonas and David.
\author{Mahault Albarracin \and Jonas Hallgren \and David Hyland}
\date{Draft, \today}

\begin{document}
\maketitle

\begin{abstract}
\todo{Write once P1 is in.}
Pluralist agents, who hold rival theories side by side, can track a paradigm shift
that monist agents miss \citep{hallgren2026Pluralist}. Holding a rival is not the same
as testing it. We ask when a pluralist community stops behaving pluralistically and
locks in. We treat lock-in as the collapse of precision on disconfirming evidence,
through two channels: the agent's own observations, and what its neighbours report.
\todo{Result: does pluralism protect against collapse of one channel, both, or neither?}
\end{abstract}

\section{Introduction}
% Feyerabend's pluralism \citep{feyerabend1975AgainstMethod} and the IWAI result.
% The gap: pluralism in belief vs pluralism in practice. A community can hold a rival
% and still never let evidence for it through.
% Our account: two channels, one currency (precision). Conviction and resource limits
% close the world channel; trust governs the social channel.
% Contribution list once results are in.
\todo{Draft after P1.}

\section{Background}
\subsection{Pluralist and monist agents}
% Summary of the IWAI model: K Gaussian theories in information form, support weights,
% support sharing vs belief sharing \citep{hallgren2026Pluralist}.
\subsection{Precision as the gate on disconfirming evidence}
% Precision-weighting in active inference \citep{feldmanFriston2010Attention,parrFriston2019Attention}.
% Motivated cost of changing one's mind \citep{hylandAlbarracin2025MindChange}.
\subsection{Network epistemology}
% \citep{zollman2010TransientDiversity,oconnorWeatherall2019Misinformation,albarracin2022EpistemicCommunities}

\section{Model}
\subsection{Base: the pluralist model}
% src/pluralist/model.py (port of David's epistemic-pluralism code).
\subsection{World channel}
% Endogenous attention on the discriminating channel:
% a_i = a_0 * w(conviction_i) * g(r_i). Conviction gate and resource gate.
\subsection{Social channel}
% Trust learned from neighbours' predictive density, Gamma-conjugate update.
\subsection{Belief utility and the motivated cost of mind-change}
% Gate on negative moves in support for the valued theory (P2).
\subsection{Support forgetting}
% The IWAI model never forgets support. We add a decay rho_L and report results
% with and without it, so a pluralist advantage is not an artefact of saturation.

\section{Experiments}
\subsection{Protocol}
% Shift protocol: mass sign flips at round 80; full observers plus restricted agents.
\subsection{P1: pluralist vs monist under two-channel collapse}
\todo{P1 results: experiments/results/P1/}
\subsection{P2: belief utility}
\todo{P2 results}

\section{Discussion}
% Pluralism in belief vs in practice. Willful vs deep ignorance
% \citep{freedman2024KnowingBetter} as the bridge to the motivated-ignorance follow-up
% \citep{williams2021MotivatedIgnorance}. Trust split and selective sharing as the
% second follow-up \citep{mayer1995Trust,hagenbachSaucet2025MotivatedSkepticism}.

\bibliographystyle{plainnat}
\bibliography{refs}
\end{document}
